\documentclass[10pt]{extarticle}
\usepackage[a4paper,margin=1.5cm]{geometry}
\usepackage[T1]{fontenc}
\usepackage{amsmath,amssymb}
\usepackage{graphicx}
\usepackage[export]{adjustbox}
\usepackage{multicol}
\usepackage{enumitem}
\usepackage{booktabs,array,tabularx}
\setlength{\columnsep}{0.7cm}
\setlength{\parindent}{0pt}
\setlength{\parskip}{3pt}
\pagestyle{empty}
\makeatletter
% figure registry (figdefs.tex is generated by the extract step)
\newcommand{\qfigdef}[4]{\@namedef{qf@file@#1}{#2}\@namedef{qf@w@#1}{#3}\@namedef{qf@wide@#1}{#4}}
\newcommand{\qf@check}[1]{\@ifundefined{qf@file@#1}{\PackageError{score2}{Unknown figure id #1}{}}{}}
\newcommand{\qf@img}[4][]{\includegraphics[width=#2\dimexpr\@nameuse{qf@w@#4}pt\relax,max width=#3,max height=0.45\textheight,#1]{\@nameuse{qf@file@#4}}}
% Full-width questions: finalize wraps every question that contains a wide figure in \qfullbegin...\qfullend,
% so the WHOLE question (stem, figures, options) is set across the page in reading order, then the two columns
% resume. The question's head - its number, stem and first figure (\qfullsplit marks the end of the head) -
% is kept on one page, so a number is never left at a page foot with its figure overleaf; the rest may break
% across pages normally, so no big gaps are left. Without a \qfullsplit, the whole question is kept together
% if it is under 45% of a page. Never write these yourself.
\newif\ifqf@full
\newif\ifqf@boxed
\newsavebox{\qf@box}
\newcommand{\qf@open}{\global\qf@fulltrue\par\global\qf@boxedtrue
  \setbox\qf@box\vbox\bgroup\hsize\textwidth\linewidth\textwidth}
\newcommand{\qf@emit}[1]{\par\egroup\global\qf@boxedfalse
  \ifdim\dimexpr\ht\qf@box+\dp\qf@box\relax<#1\textheight\noindent\usebox{\qf@box}\par\else\unvbox\qf@box\fi}
\newcommand{\qfullsplit}{\ifqf@boxed\qf@emit{0.85}\fi}
\newcommand{\qf@close}{\ifqf@boxed\qf@emit{0.45}\fi\global\qf@fullfalse}
\newcommand{\qfullbegin}{\end{multicols}\qf@open}
\newcommand{\qfullend}{\qf@close\begin{multicols}{2}}
\newcommand{\qfullstart}{\qf@open}
\newcommand{\qfulldone}{\qf@close}
% \qfig[scale]{id}: block figure, centred, directly where it is written (size set by the tool).
% Fallback if a wide figure is not inside a full-width question: break out of the columns in place - no floats.
\newcommand{\qfig}[2][1]{\qf@check{#2}%
  \ifqf@full\qf@blk{#1}{#2}\else
  \ifnum\@nameuse{qf@wide@#2}=1 \expandafter\@firstoftwo\else\expandafter\@secondoftwo\fi
  {\end{multicols}\begin{center}\qf@img{#1}{\textwidth}{#2}\end{center}\begin{multicols}{2}}%
  {\qf@blk{#1}{#2}}\fi}
\newcommand{\qf@blk}[2]{\par\smallskip{\centering\qf@img{#1}{\linewidth}{#2}\par}\smallskip}
% \qfigi[scale]{id}: inline figure (inside option lists / tables / minipages); never breaks columns.
\newcommand{\qfigi}[2][1]{\qf@check{#2}\qf@img[valign=c]{#1}{\linewidth}{#2}}
\makeatother
% structure macros - lint reads these to check completeness and figure placement
\newcommand{\papersec}[1]{\par\bigskip{\centering\Large\bfseries PAPER-#1\par}\medskip}
\newcommand{\partsec}[2]{\par\medskip\noindent{\large\bfseries PART-#1: #2}\par\smallskip}
\newcommand{\sect}[1]{\par\smallskip\noindent\textbf{#1}\par}
\newcommand{\qn}[1]{\par\smallskip\noindent\textbf{#1.}~}
\newcommand{\paraq}[2]{\par\smallskip\noindent\textbf{Paragraph for Questions #1 and #2}\par}
\newenvironment{qoptions}{\begin{itemize}[leftmargin=1.8em,itemsep=1pt,topsep=2pt,label={}]}{\end{itemize}}
\newcommand{\opt}[2]{\item[(#1)] #2}
% figure registry (generated by score2_tools.py finalize)
\qfigdef{p1_phy_s1_q02_1}{Internal Test - Test 01 (05 Aug 2020)-img/p1_phy_s1_q02_1.png}{146.6}{0}
\qfigdef{p1_phy_s1_q03_1}{Internal Test - Test 01 (05 Aug 2020)-img/p1_phy_s1_q03_1.png}{270.4}{1}
\qfigdef{p1_phy_s1_pre_1}{Internal Test - Test 01 (05 Aug 2020)-img/p1_phy_s1_pre_1.png}{134.0}{0}
\qfigdef{p1_phy_s2_q05_1}{Internal Test - Test 01 (05 Aug 2020)-img/p1_phy_s2_q05_1.png}{174.4}{0}
\qfigdef{p1_phy_s2_pre_1}{Internal Test - Test 01 (05 Aug 2020)-img/p1_phy_s2_pre_1.png}{75.7}{0}
\qfigdef{p1_phy_s2_pre_2}{Internal Test - Test 01 (05 Aug 2020)-img/p1_phy_s2_pre_2.png}{243.3}{0}
\qfigdef{p1_phy_s2_pre_3}{Internal Test - Test 01 (05 Aug 2020)-img/p1_phy_s2_pre_3.png}{146.4}{0}
\qfigdef{p1_phy_s3_q11_1}{Internal Test - Test 01 (05 Aug 2020)-img/p1_phy_s3_q11_1.png}{322.3}{1}
\qfigdef{p1_phy_s3_q12_1}{Internal Test - Test 01 (05 Aug 2020)-img/p1_phy_s3_q12_1.png}{219.1}{0}
\qfigdef{p1_phy_s3_q13_1}{Internal Test - Test 01 (05 Aug 2020)-img/p1_phy_s3_q13_1.png}{226.0}{0}
\qfigdef{p1_chem_s1_q01_1}{Internal Test - Test 01 (05 Aug 2020)-img/p1_chem_s1_q01_1.png}{354.5}{1}
\qfigdef{p1_chem_s1_q01_A}{Internal Test - Test 01 (05 Aug 2020)-img/p1_chem_s1_q01_A.png}{263.1}{1}
\qfigdef{p1_chem_s1_q01_B}{Internal Test - Test 01 (05 Aug 2020)-img/p1_chem_s1_q01_B.png}{277.7}{1}
\qfigdef{p1_chem_s1_q01_C}{Internal Test - Test 01 (05 Aug 2020)-img/p1_chem_s1_q01_C.png}{103.2}{0}
\qfigdef{p1_chem_s1_q01_D}{Internal Test - Test 01 (05 Aug 2020)-img/p1_chem_s1_q01_D.png}{99.9}{0}
\qfigdef{p1_chem_s1_pre_1}{Internal Test - Test 01 (05 Aug 2020)-img/p1_chem_s1_pre_1.png}{115.0}{0}
\qfigdef{p1_chem_s1_pre_A}{Internal Test - Test 01 (05 Aug 2020)-img/p1_chem_s1_pre_A.png}{92.1}{0}
\qfigdef{p1_chem_s1_pre_B}{Internal Test - Test 01 (05 Aug 2020)-img/p1_chem_s1_pre_B.png}{91.5}{0}
\qfigdef{p1_chem_s1_pre_C}{Internal Test - Test 01 (05 Aug 2020)-img/p1_chem_s1_pre_C.png}{92.1}{0}
\qfigdef{p1_chem_s1_pre_D}{Internal Test - Test 01 (05 Aug 2020)-img/p1_chem_s1_pre_D.png}{87.3}{0}
\qfigdef{p1_chem_s1_pre_6}{Internal Test - Test 01 (05 Aug 2020)-img/p1_chem_s1_pre_6.png}{137.9}{0}
\qfigdef{p1_chem_s1_pre_7A}{Internal Test - Test 01 (05 Aug 2020)-img/p1_chem_s1_pre_7A.png}{51.1}{0}
\qfigdef{p1_chem_s1_pre_8B}{Internal Test - Test 01 (05 Aug 2020)-img/p1_chem_s1_pre_8B.png}{33.4}{0}
\qfigdef{p1_chem_s1_pre_9C}{Internal Test - Test 01 (05 Aug 2020)-img/p1_chem_s1_pre_9C.png}{87.1}{0}
\qfigdef{p1_chem_s1_pre_10D}{Internal Test - Test 01 (05 Aug 2020)-img/p1_chem_s1_pre_10D.png}{38.8}{0}
\qfigdef{p1_chem_s1_pre_11A}{Internal Test - Test 01 (05 Aug 2020)-img/p1_chem_s1_pre_11A.png}{198.1}{0}
\qfigdef{p1_chem_s1_q04_A}{Internal Test - Test 01 (05 Aug 2020)-img/p1_chem_s1_q04_A.png}{111.9}{0}
\qfigdef{p1_chem_s1_q04_B}{Internal Test - Test 01 (05 Aug 2020)-img/p1_chem_s1_q04_B.png}{107.1}{0}
\qfigdef{p1_chem_s1_q04_C}{Internal Test - Test 01 (05 Aug 2020)-img/p1_chem_s1_q04_C.png}{112.2}{0}
\qfigdef{p1_chem_s1_q04_D}{Internal Test - Test 01 (05 Aug 2020)-img/p1_chem_s1_q04_D.png}{94.5}{0}
\qfigdef{p1_chem_s2_q05_B}{Internal Test - Test 01 (05 Aug 2020)-img/p1_chem_s2_q05_B.png}{148.5}{0}
\qfigdef{p1_chem_s2_q05_C}{Internal Test - Test 01 (05 Aug 2020)-img/p1_chem_s2_q05_C.png}{33.6}{0}
\qfigdef{p1_chem_s2_pre_1}{Internal Test - Test 01 (05 Aug 2020)-img/p1_chem_s2_pre_1.png}{269.9}{1}
\qfigdef{p1_chem_s2_pre_2}{Internal Test - Test 01 (05 Aug 2020)-img/p1_chem_s2_pre_2.png}{60.9}{0}
\qfigdef{p1_chem_s2_pre_B}{Internal Test - Test 01 (05 Aug 2020)-img/p1_chem_s2_pre_B.png}{71.8}{0}
\qfigdef{p1_chem_s2_pre_4}{Internal Test - Test 01 (05 Aug 2020)-img/p1_chem_s2_pre_4.png}{77.0}{0}
\qfigdef{p1_chem_s2_q07_A}{Internal Test - Test 01 (05 Aug 2020)-img/p1_chem_s2_q07_A.png}{125.9}{0}
\qfigdef{p1_chem_s2_q07_B}{Internal Test - Test 01 (05 Aug 2020)-img/p1_chem_s2_q07_B.png}{125.9}{0}
\qfigdef{p1_chem_s2_q07_C}{Internal Test - Test 01 (05 Aug 2020)-img/p1_chem_s2_q07_C.png}{135.7}{0}
\qfigdef{p1_chem_s2_q07_D}{Internal Test - Test 01 (05 Aug 2020)-img/p1_chem_s2_q07_D.png}{127.2}{0}
\qfigdef{p1_chem_s2_q08_A}{Internal Test - Test 01 (05 Aug 2020)-img/p1_chem_s2_q08_A.png}{205.3}{0}
\qfigdef{p1_chem_s2_q08_B}{Internal Test - Test 01 (05 Aug 2020)-img/p1_chem_s2_q08_B.png}{180.2}{0}
\qfigdef{p1_chem_s2_pre_15C}{Internal Test - Test 01 (05 Aug 2020)-img/p1_chem_s2_pre_15C.png}{178.7}{0}
\qfigdef{p1_chem_s2_q08_D}{Internal Test - Test 01 (05 Aug 2020)-img/p1_chem_s2_q08_D.png}{197.9}{0}
\qfigdef{p1_chem_s2_q09_A}{Internal Test - Test 01 (05 Aug 2020)-img/p1_chem_s2_q09_A.png}{176.3}{0}
\qfigdef{p1_chem_s2_q09_B}{Internal Test - Test 01 (05 Aug 2020)-img/p1_chem_s2_q09_B.png}{348.5}{1}
\qfigdef{p1_chem_s2_q09_C}{Internal Test - Test 01 (05 Aug 2020)-img/p1_chem_s2_q09_C.png}{174.5}{0}
\qfigdef{p1_chem_s2_pre_16D}{Internal Test - Test 01 (05 Aug 2020)-img/p1_chem_s2_pre_16D.png}{198.7}{0}
\qfigdef{p1_chem_s2_pre_17}{Internal Test - Test 01 (05 Aug 2020)-img/p1_chem_s2_pre_17.png}{118.5}{0}
\qfigdef{p1_chem_s3_q11_A}{Internal Test - Test 01 (05 Aug 2020)-img/p1_chem_s3_q11_A.png}{168.0}{0}
\qfigdef{p1_chem_s3_q11_B}{Internal Test - Test 01 (05 Aug 2020)-img/p1_chem_s3_q11_B.png}{126.4}{0}
\qfigdef{p1_chem_s3_q11_C}{Internal Test - Test 01 (05 Aug 2020)-img/p1_chem_s3_q11_C.png}{167.5}{0}
\qfigdef{p1_chem_s3_q11_D}{Internal Test - Test 01 (05 Aug 2020)-img/p1_chem_s3_q11_D.png}{192.6}{0}
\qfigdef{p1_chem_s3_q11_F}{Internal Test - Test 01 (05 Aug 2020)-img/p1_chem_s3_q11_F.png}{209.3}{0}
\qfigdef{p1_chem_s3_pre_D}{Internal Test - Test 01 (05 Aug 2020)-img/p1_chem_s3_pre_D.png}{51.1}{0}
\qfigdef{p1_chem_s3_pre_2}{Internal Test - Test 01 (05 Aug 2020)-img/p1_chem_s3_pre_2.png}{98.2}{0}
\qfigdef{p1_chem_s3_pre_3}{Internal Test - Test 01 (05 Aug 2020)-img/p1_chem_s3_pre_3.png}{92.7}{0}
\qfigdef{p1_chem_s3_pre_4}{Internal Test - Test 01 (05 Aug 2020)-img/p1_chem_s3_pre_4.png}{117.8}{0}
\qfigdef{p1_chem_s3_pre_5}{Internal Test - Test 01 (05 Aug 2020)-img/p1_chem_s3_pre_5.png}{72.6}{0}
\qfigdef{p1_chem_s3_pre_6}{Internal Test - Test 01 (05 Aug 2020)-img/p1_chem_s3_pre_6.png}{89.6}{0}
\qfigdef{p1_chem_s3_pre_7}{Internal Test - Test 01 (05 Aug 2020)-img/p1_chem_s3_pre_7.png}{66.4}{0}
\qfigdef{p1_chem_s3_pre_8}{Internal Test - Test 01 (05 Aug 2020)-img/p1_chem_s3_pre_8.png}{33.4}{0}
\qfigdef{p1_chem_s3_pre_9}{Internal Test - Test 01 (05 Aug 2020)-img/p1_chem_s3_pre_9.png}{81.8}{0}
\qfigdef{p1_chem_s3_pre_10}{Internal Test - Test 01 (05 Aug 2020)-img/p1_chem_s3_pre_10.png}{75.7}{0}
\qfigdef{p1_chem_s3_pre_11}{Internal Test - Test 01 (05 Aug 2020)-img/p1_chem_s3_pre_11.png}{90.1}{0}
\qfigdef{p1_chem_s3_pre_12}{Internal Test - Test 01 (05 Aug 2020)-img/p1_chem_s3_pre_12.png}{130.7}{0}
\raggedcolumns
\makeatletter\@beginparpenalty=10000\makeatother
\newcolumntype{L}[1]{>{\raggedright\arraybackslash}p{#1}}
\begin{document}
\begin{center}{\Large\bfseries Internal Test --- Test 01 (05 Aug 2020)}\end{center}
\begin{multicols}{2}
\raggedcolumns
\papersec{1}
\partsec{1}{PHYSICS}
\sect{SECTION-I(i)}
\qn{1} Under the influence of the coulomb field of charge $+Q$ which is fixed, a charge $-q$ is moving around it in an elliptical orbit. Find out correct statements
\begin{qoptions}
\opt{A}{The angular momentum of the charge $-q$ is constant.}
\opt{B}{The linear momentum of the charge $-q$ is constant.}
\opt{C}{The angular velocity of the charge $-q$ is constant.}
\opt{D}{The linear speed of the charge $-q$ is constant.}
\end{qoptions}
\qn{2} Find the magnitude of dipole moment of the following system.\par
\qfig{p1_phy_s1_q02_1}
\begin{qoptions}
\opt{A}{$\sqrt{21}\,aq$}
\opt{B}{Zero}
\opt{C}{$4aq$}
\opt{D}{$10aq$}
\end{qoptions}
\qfullbegin
\qn{3} The potential difference between A and B in the following situation is? $(V_A-V_B)$\par
\qfig{p1_phy_s1_q03_1}
\qfullsplit
\begin{qoptions}
\opt{A}{$12$ V}
\opt{B}{$7$ V}
\opt{C}{$20$ V}
\opt{D}{$14$ V}
\end{qoptions}
\qfullend
\qn{4} In the circuit shown in figure the heat produced in the $5\,\Omega$ resistor due to the current flowing through it is $10$ calories per second.\par
\qfig{p1_phy_s1_pre_1}
The heat generated in the $4\,\Omega$ resistor is :
\begin{qoptions}
\opt{A}{$1$ cal/s}
\opt{B}{$2$ cal/s}
\opt{C}{$3$ cal/s}
\opt{D}{$4$ cal/s}
\end{qoptions}
\sect{SECTION-I(ii)}
\qn{5} A current enters from A and leaves from B through the network shown. Select the correct alternatives\par
\qfig{p1_phy_s2_q05_1}
\begin{qoptions}
\opt{A}{$R_{AB}=\dfrac{3R}{4}$}
\opt{B}{$R_{AB}=\dfrac{5R}{6}$}
\opt{C}{Points P, Q, T and S are at same potential}
\opt{D}{Points C and D are at same potential}
\end{qoptions}
\qn{6} Consider gauss law $\oint \vec E\cdot d\vec A=\dfrac{q}{\epsilon_0}$, which of the following is not true?
\begin{qoptions}
\opt{A}{$\vec E$ must be the electric field due to enclosed charge only}
\opt{B}{If net charge inside the gaussian surface $=0$, $\vec E$ must be zero everywhere over the gaussian surface}
\opt{C}{If the only charge inside the gaussian surface is an electric dipole, then the integral is zero}
\opt{D}{$\vec E$ is paralled to $d\vec A$ everywhere over the gaussian surface.}
\end{qoptions}
\qn{7} A dipole with electric dipole moment of $\vec p$ is located at a distance $\vec r$ from a long thread charged uniformly with a linear charge density $\lambda$ as shown.\par
\qfig{p1_phy_s2_pre_1}
\begin{qoptions}
\opt{A}{The force on dipole is non zero}
\opt{B}{The force acting on dipole if the vector $\vec p$ is oriented along the thread is $\vec F=0$}
\opt{C}{The force acting on dipole if the vector $\vec p$ is oriented along the $\vec r$ is $\vec F=-\dfrac{2\lambda\vec P}{\pi\epsilon_0 r^2}$.}
\opt{D}{The force acting on the dipole if the vector $\vec p$ is oriented at right angles to thread is $|\vec F|=\dfrac{\lambda p}{2\pi\epsilon_0 r^2}$}
\end{qoptions}
\qn{8} A point charge $q$, is placed at a distance $r$ from the centre of an uncharged aluminium sphere, of radius $R\ (r>R)$. Select the correct alternatives.\par
\qfig[0.85]{p1_phy_s2_pre_2}
\begin{qoptions}
\opt{A}{When the switch is open, the potential at ``O'' is $\dfrac{q}{4\pi\epsilon_0 r}$}
\opt{B}{When the switch is closed, the potential at ``O'' is zero}
\opt{C}{A charge $-q$ will be present on sphere after closing switch.}
\opt{D}{When switch is closed, sphere acquires a net charge $-\dfrac{qR}{r}$}
\end{qoptions}
\qn{9} A spherical body of radius $R$ consists of a fluid of constant density and is in equilibrium under its own gravity. If $P(r)$ is pressure at $r$ ($r<R$), then the correct options is (are)
\begin{qoptions}
\opt{A}{$P(r=0)=0$}
\opt{B}{$\dfrac{P\left(r=\frac{3R}{2}\right)}{P\left(r=\frac{2R}{3}\right)}=\dfrac{63}{80}$}
\opt{C}{$\dfrac{P\left(r=\frac{3R}{5}\right)}{P\left(r=\frac{2R}{5}\right)}=\dfrac{16}{21}$}
\opt{D}{$\dfrac{P\left(r=\frac{R}{2}\right)}{P\left(r=\frac{R}{3}\right)}=\dfrac{20}{27}$}
\end{qoptions}
\qn{10} Figure shows a conical conducting wire connected to a source of emf ``e''. Let $E, V_d, i$ represent the electric field, drift velocity, and current at a cross section of the wire. As one moves from end A to end B of the wire\par
\qfig[0.8]{p1_phy_s2_pre_3}
\begin{qoptions}
\opt{A}{$E$ increases}
\opt{B}{$E$ decreases}
\opt{C}{$V_d$ decreases}
\opt{D}{$i$ remains same}
\end{qoptions}
\qfullbegin
\sect{SECTION-II}
\qn{11} Consider the infinite ladder circuit shown. What is the value of current $I$ (in A) shown in one of the $6\,\Omega$ resistors?\par
\qfig{p1_phy_s3_q11_1}
\qfullend
\qn{12} Consider the circuit shown in figure the current (in ampere) through $2\,\Omega$ resistor is\par
\qfig[0.75]{p1_phy_s3_q12_1}
\qn{13} Find the resistance between A and B of the circuit shown in figure.\par
\qfig{p1_phy_s3_q13_1}
\qn{14} $1$m long metallic wire is broken into two unequal parts P and Q. P of the wire is uniformly extended into another wire R. Length of R is twice the length of P and the resistance of R is equal to that of Q. Find the ratio of the length of Q and P.
\qn{15} A bullet is fired vertically upward with velocity ``$V$'' from the surface of a spherical planet. When it reaches its maximum height, its acceleration due to the planet's gravity is $\dfrac1{16}$th of its value at the surface of the planet. If the escape velocity from the planet surface is $V_{esc}=V\sqrt N$, then the value of $N$ is (ignore atmosphere).
\qn{16} A solid sphere of radius ``$R$'' has a charge $Q$ distributed in its volume with a charge density $\rho=Kr^a$, where $K$ and $a$ are constants and $r$ is the distance from its centre. If the electric field at $r=\dfrac R2$ is $\dfrac1{16}$ times that at $r=R$, find the value of $a$.
\qn{17} An electric field is given by $\vec E=(y\hat i+x\hat j)\,$N/C. What is the work done in moving a $1$C charge from $\vec r_A=(2\hat i+2\hat j)$ to $\vec r_B=(4\hat i+\hat j)$?
\qn{18} Two point objects of masses $m$ and $4m$ are at rest at an infinite seperation. A slight push is given to one of them and they move towards each other under mutual gravitational attraction. If $G$ is the universal gravitational constant. Then at a seperation ``$r$'' total K.E of objects is given by $\dfrac{xGm^2}{r}$ find $x$?
\qfullbegin
\partsec{2}{CHEMISTRY}
\sect{SECTION-I(i)}
\qn{1} Identify product of the given reactions :\par
\qfig{p1_chem_s1_q01_1}
\qfullsplit
Product (1) and (2) respectively
\begin{qoptions}
\opt{A}{\qfigi{p1_chem_s1_q01_A}}
\opt{B}{\qfigi{p1_chem_s1_q01_B}}
\opt{C}{\qfigi{p1_chem_s1_q01_C} single product is obtained in both the reactions}
\opt{D}{\qfigi{p1_chem_s1_q01_D} single product is obtained in both the reactions}
\end{qoptions}
\qfullend
\qn{2} \qfig{p1_chem_s1_pre_1}
The product `P' is
\begin{qoptions}
\opt{A}{\qfigi{p1_chem_s1_pre_A}}
\opt{B}{\qfigi{p1_chem_s1_pre_B}}
\opt{C}{\qfigi{p1_chem_s1_pre_C}}
\opt{D}{\qfigi{p1_chem_s1_pre_D}}
\end{qoptions}
\qn{3} \qfig{p1_chem_s1_pre_6}
\begin{qoptions}
\opt{A}{\qfigi{p1_chem_s1_pre_7A}}
\opt{B}{\qfigi{p1_chem_s1_pre_8B}}
\opt{C}{\qfigi{p1_chem_s1_pre_9C}}
\opt{D}{\qfigi{p1_chem_s1_pre_10D}}
\end{qoptions}
\qn{4} Identify the major product in the following reaction\par
\qfig{p1_chem_s1_pre_11A}
\begin{qoptions}
\opt{A}{\qfigi{p1_chem_s1_q04_A}}
\opt{B}{\qfigi{p1_chem_s1_q04_B}}
\opt{C}{\qfigi{p1_chem_s1_q04_C}}
\opt{D}{\qfigi{p1_chem_s1_q04_D}}
\end{qoptions}
\sect{SECTION-I(ii)}
\qn{5} Identify the correctly drawn conclusions about alcohols and phenols.
\begin{qoptions}
\opt{A}{Alcohols which cannot produce any specific colour in Victor meyer test will give turbidity instantaneously in Lucas test}
\opt{B}{\qfigi{p1_chem_s2_q05_B} can be distinguished through reaction with $\mathrm{HIO_4}$}
\opt{C}{$\mathrm{CH_3-C(OH){=}CH-CH_3}$ and \qfigi{p1_chem_s2_q05_C} can be distinguished through reaction with neutral $\mathrm{FeCl_3}$}
\opt{D}{Ethyl alcohol and formaldehyde can be distinguished through reaction with NaOI}
\end{qoptions}
\qfullbegin
\qn{6} Identify the correct statements in the below reaction scheme.\par
\qfig{p1_chem_s2_pre_1}
\qfullsplit
\begin{qoptions}
\opt{A}{A $=$ \qfigi{p1_chem_s2_pre_2}}
\opt{B}{D $=$ \qfigi{p1_chem_s2_pre_B}}
\opt{C}{E $=$ \qfigi{p1_chem_s2_pre_4}}
\opt{D}{F $=$ Bakelite}
\end{qoptions}
\qfullend
\qn{7} Select the incorrect reactions.
\begin{qoptions}
\opt{A}{\qfigi{p1_chem_s2_q07_A}}
\opt{B}{\qfigi{p1_chem_s2_q07_B}}
\opt{C}{\qfigi{p1_chem_s2_q07_C}}
\opt{D}{\qfigi{p1_chem_s2_q07_D}}
\end{qoptions}
\qn{8} Which of the following reactions is/are correctly matched :
\begin{qoptions}
\opt{A}{\qfigi{p1_chem_s2_q08_A}}
\opt{B}{\qfigi{p1_chem_s2_q08_B}}
\opt{C}{\qfigi{p1_chem_s2_pre_15C}}
\opt{D}{\qfigi{p1_chem_s2_q08_D}}
\end{qoptions}
\qfullbegin
\qn{9} Which of the following reaction involve hydride shift ?
\begin{qoptions}
\opt{A}{\qfigi{p1_chem_s2_q09_A}}
\opt{B}{\qfigi{p1_chem_s2_q09_B}}
\opt{C}{\qfigi{p1_chem_s2_q09_C}}
\opt{D}{\qfigi{p1_chem_s2_pre_16D}}
\end{qoptions}
\qfullend
\qn{10} Which of the following changes are involved in the formation of most stable product X under given condition.\par
\qfig{p1_chem_s2_pre_17}
\begin{qoptions}
\opt{A}{Delocalisation of charge}
\opt{B}{1,2 shift of Ethyl group}
\opt{C}{Aromatisation of ring in which ethyl substituent is present}
\opt{D}{Aromatisation of ring in which ethyl substituent is absent}
\end{qoptions}
\sect{SECTION-II}
\qn{11} How many of the conclusions drawn about in the each of the following reaction takes place under suitable conditions are correct
\begin{qoptions}
\opt{A}{Both H$^+$ and OH$^-$ can catalyse the reaction\par\smallskip\qfigi{p1_chem_s3_q11_A}}
\opt{B}{Both H$^+$ and OH$^-$ can catalyse the reaction\par\smallskip\qfigi{p1_chem_s3_q11_B}}
\opt{C}{Only H$^+$ can catalyse the reaction\par\smallskip\qfigi{p1_chem_s3_q11_C}}
\opt{D}{Only H$^+$ can catalyse the reaction\par\smallskip\qfigi{p1_chem_s3_q11_D}}
\opt{E}{In Perkins condensation the internal transfer of the acetyl group takes place from the carboxyl oxygen atom to the alkoxy oxygen atom via cyclic six membered intermediate.}
\opt{F}{\qfigi{p1_chem_s3_q11_F}. The degree of unsaturation of Y is 5.}
\end{qoptions}
\qn{12} How many of the following compounds are oxidized by $\mathrm{HIO_4}$.
\begin{qoptions}
\opt{A}{$\mathrm{CH_3{-}CO{-}CHO}$}
\opt{B}{$\mathrm{CH_3{-}CO{-}CO{-}CH_3}$}
\opt{C}{$\mathrm{CH_3{-}CO{-}CH(NH_2){-}CH_3}$}
\opt{D}{\qfigi{p1_chem_s3_pre_D}}
\opt{E}{$\mathrm{HOCH_2{-}CH_2{-}OCH_3}$}
\opt{F}{$\mathrm{HOCH_2{-}CO{-}CH_2OH}$}
\opt{G}{$\mathrm{HOCH_2{-}CH_2{-}CH_2OH}$}
\opt{H}{\qfigi{p1_chem_s3_pre_2}}
\opt{I}{\qfigi{p1_chem_s3_pre_3}}
\end{qoptions}
\qn{13} In the below reaction of Gem Dihalide formation, pick how many statements are correct.\par
{\centering $\mathrm{CH_3{-}CHO}\ \xrightarrow{\mathrm{PCl_5}}\ \mathrm{CH_3{-}CH(Cl)Cl}$\par}
\begin{qoptions}
\opt{A}{First step involves attack of $\mathrm{PCl_4^+}$ at carbonyl oxygen.}
\opt{B}{The second chloride ion comes from $\mathrm{PCl_6^-}$}
\opt{C}{Carbocation formation at carbonyl carbon takes place after first step.}
\opt{D}{Internal nucleophilic attack (SNi) is also possible to bring second chloride ion.}
\end{qoptions}
\qn{14} Pick how many options shown below will work to prepare 1, 1, 1, 2, 2-pentadeutero-3, 3-dimethylpentane from 3, 3-Dimethyl-2-pentanone.
\begin{qoptions}
\opt{A}{(i) NaOD in $\mathrm{D_2O}$ (excess), (ii) $\mathrm{LiAlD_4}$ in ether, (iii) $\mathrm{D_2O}$ workup.}
\opt{B}{(i) $\mathrm{N_2D_4}$/KOD, (ii) $\mathrm{NaOD/D_2O}$ excess}
\opt{C}{(i) Zn/DCl in $\mathrm{CH_3COOD}$, (ii) $\mathrm{LiAlD_4}$ in ether, (iii) $\mathrm{D_2O}$ workup.}
\opt{D}{(i) NaOD in $\mathrm{D_2O}$ (excess), (ii) $\mathrm{CH_2(SH){-}CH_2(SH)}$ $+$ $\mathrm{BF_3}$, (iii) $\mathrm{D_2}$ $+$ Raney Ni}
\end{qoptions}
\qn{15} Number of moles of HI consumed per mole of below compound is/are\par
\qfig{p1_chem_s3_pre_4}
\qn{16} The total number of possible aldol products including stereoisomers.\par
{\centering $\mathrm{Ph{-}CHO}\ +\ \mathrm{CH_3{-}CHO}\ \xrightarrow[5^\circ\mathrm{C}]{\mathrm{NaOH/H_2O}}$\par}
\qn{17} The number compounds which will give the Tollen's test.
\begin{qoptions}
\opt{i}{$\mathrm{D{-}CO{-}D}$}
\opt{ii}{$\mathrm{CH_3{-}CO{-}D}$}
\opt{iii}{\qfigi{p1_chem_s3_pre_5}}
\opt{iv}{$\mathrm{D{-}CO{-}OD}$}
\opt{v}{$\mathrm{CH_3{-}CO{-}CH_3}$}
\opt{vi}{$\mathrm{CH_3{-}C{\equiv}C{-}CH_3}$}
\opt{vii}{$\mathrm{CH_3{-}C{\equiv}CH}$}
\opt{viii}{$\mathrm{CH_3{-}CO{-}CH(OH){-}CH_3}$}
\opt{ix}{\qfigi{p1_chem_s3_pre_6}}
\opt{x}{$\mathrm{CH_3{-}CO{-}OH}$}
\opt{xi}{\qfigi{p1_chem_s3_pre_7}}
\end{qoptions}
\qn{18} How many of the following compounds form stable hydrates.
\begin{qoptions}
\opt{i}{HCHO}
\opt{ii}{\qfigi{p1_chem_s3_pre_8}}
\opt{iii}{$\mathrm{CCl_3CHO}$}
\opt{iv}{\qfigi{p1_chem_s3_pre_9}}
\opt{v}{$\mathrm{CH_3{-}CO{-}CH_3}$}
\opt{vi}{Glyoxal}
\opt{vii}{$\mathrm{Ph{-}CO{-}CO{-}CO{-}Ph}$}
\opt{viii}{\qfigi{p1_chem_s3_pre_10}}
\opt{ix}{\qfigi{p1_chem_s3_pre_11}}
\opt{x}{\qfigi{p1_chem_s3_pre_12}}
\end{qoptions}
\partsec{3}{MATHEMATICS}
\sect{SECTION-I(i)}
\qn{1} If $f(x)+f(y)=f\left(\dfrac{x+y}{1-xy}\right)$ for all $x,y\in\mathbb R$ and $xy\ne1$, and $\displaystyle\lim_{x\to0}\dfrac{f(x)}{x}=2$, the value of $f\left(\dfrac1{\sqrt3}\right)$ is
\begin{qoptions}
\opt{A}{$1$}
\opt{B}{$\dfrac13$}
\opt{C}{$\dfrac\pi3$}
\opt{D}{$\dfrac\pi2$}
\end{qoptions}
\qn{2} In terms of a parameter $t$ and a constant $k$, a curve is specified as follows:\par
$x=t+\dfrac1t,\ y=t^k+\dfrac1{t^k}$\par
If $\left(\dfrac{dy}{dx}\right)^2=\lambda\left(\dfrac{y^2-4}{x^2-4}\right)$, the value of $\lambda$ in terms of $k$ is
\begin{qoptions}
\opt{A}{$k^2$}
\opt{B}{$2k^2$}
\opt{C}{$\dfrac{k^2}{2}$}
\opt{D}{None of these}
\end{qoptions}
\qn{3} Let $f(x)$ and $g(x)$ be two functions such that $f(x)=x^2g(x)$, $\displaystyle\lim_{x\to0}g(x)\csc x=\dfrac{f'''(0)}{\lambda}$.\par
The value of $\lambda$ is
\begin{qoptions}
\opt{A}{$2$}
\opt{B}{$3$}
\opt{C}{$4$}
\opt{D}{$6$}
\end{qoptions}
\qn{4} Let $f(x)$ be a funciton defined as $f(x)=\begin{cases}ae^x\sin x, & x<0\\ x^2+b, & x\ge0\end{cases}$. Suppose that $f'(0)$ exists.\par
The value of $a+b+f'(0)$ is
\begin{qoptions}
\opt{A}{$1$}
\opt{B}{$2$}
\opt{C}{$3$}
\opt{D}{None of these}
\end{qoptions}
\sect{SECTION-I(ii)}
\qn{5} Let $f\left(\dfrac{x+y}{3}\right)=\dfrac{f(x)+f(y)+2}{3}\ \forall x,y\in\mathbb R$. Also it is given that $f'(0)=2$. Which of the following statements is/are true?
\begin{qoptions}
\opt{A}{$f'(x)=f'(0)\ \forall x\in\mathbb R$}
\opt{B}{$f(x)$ is a linear function}
\opt{C}{$f'(x)=f(x)\ \forall x\in\mathbb R$}
\opt{D}{None of these}
\end{qoptions}
\qn{6} Consider the funciton $f(x)=[x]+\sqrt{\{x\}}$. This function is where $[.]$ denotes GIF \& $\{\ \}$ denotes fractional part of $x$
\begin{qoptions}
\opt{A}{Continuous at all integer points}
\opt{B}{Discontinuous at all integer points}
\opt{C}{Differentiable at all integer points}
\opt{D}{Non-differentiable at all integer points}
\end{qoptions}
\qn{7} If $y=e^{\sqrt x}+e^{-\sqrt x}$, then $\dfrac{dy}{dx}$ is equal to
\begin{qoptions}
\opt{A}{$\dfrac{e^{\sqrt x}-e^{-\sqrt x}}{2\sqrt x}$}
\opt{B}{$\dfrac{e^{\sqrt x}-e^{-\sqrt x}}{2x}$}
\opt{C}{$\dfrac1{2\sqrt x}\sqrt{y^2-4}$}
\opt{D}{$\dfrac1{2\sqrt x}\sqrt{y^2+4}$}
\end{qoptions}
\qn{8} Which of the following is/are true?
\begin{qoptions}
\opt{A}{$\dfrac{dy}{dx}$ for $y=\sin^{-1}(\cos x)$, where $x\in(0,\pi)$, is $-1$}
\opt{B}{$\dfrac{dy}{dx}$ for $y=\sin^{-1}(\cos x)$, where $x\in(\pi,2\pi)$ is $1$}
\opt{C}{$\dfrac{dy}{dx}$ for $y=\cos^{-1}(\sin x)$, where $x\in\left(-\dfrac\pi2,\dfrac\pi2\right)$, is $-1$}
\opt{D}{$\dfrac{dy}{dx}$ for $y=\cos^{-1}(\sin x)$, where $x\in\left(\dfrac\pi2,\dfrac{3\pi}2\right)$, is $-1$}
\end{qoptions}
\qn{9} If $f(x)=\left|x^2-4x+3\right|\left[x^2\right]$ then which of the statements are correct. (where $[.]$ denotes greatest integer function).
\begin{qoptions}
\opt{A}{$f(x)$ is discontinuous at $x=1\,\&\,3$}
\opt{B}{$f(x)$ is continuous at $x=1,3$}
\opt{C}{$f'(2)$ does not exist}
\opt{D}{$f'(2)$ exist and equal to $0$.}
\end{qoptions}
\qn{10} Let $f(x)=\begin{cases}x^n\sin\left(\dfrac1{x^2}\right), & x\ne0\\ 0, & x=0\end{cases},\ (n\in I)$ then
\begin{qoptions}
\opt{A}{$\displaystyle\lim_{x\to0}f(x)$ exists for every $n>1$}
\opt{B}{$f$ is continuous at $x=0$ for $n>1$}
\opt{C}{$f$ is differentiable at $x=0$ for every $n>1$}
\opt{D}{$f$ is differentiable at $x=0$ for all values of $n$}
\end{qoptions}
\sect{SECTION-II}
\qn{11} If $y=\dfrac{x^4+x^2+1}{x^2+x+1}$ then $\dfrac{dy}{dx}=ax+b$ then the value of $a^2+b^2$ is \underline{\hspace{2cm}}
\qn{12} If $y=f(x)$ satisfies $\left(xf(x)\right)^{99}+\left(xf(x)\right)^{98}+\ldots\ldots+\left(xf(x)\right)+1=0$ then $f'(2)=$ \underline{\hspace{2cm}}
\qn{13} Number of points of non-differentiability of funciton $f(x)=\max\left\{\sin^{-1}|\sin x|,\cos^{-1}|\sin x|\right\},0<x<2\pi$ is \underline{\hspace{2cm}}
\qn{14} Let $f(x)$ be defined as $f(x)=\begin{cases}x^{\frac{k^2-3k+2}{k^2-5k+4}}\cos\left(\dfrac1x\right), & x\ne0\\ 0, & x=0\end{cases}$ then smallest positive integral value of $k$ for which $f(x)$ is differentiable at $x=0$, is \underline{\hspace{2cm}}
\qn{15} Let $f(x)=x^2-x+k-2$, $k\in R$, and the complete set of value of $k$ for which $y=\left|f\left(|x|\right)\right|$ is non-derivable at $5$ distinct point $(a,b)$ then $a-b$ is \underline{\hspace{2cm}}
\qn{16} If $y=\left(\tan^{-1}x\right)^2$ then $\left(1+x^2\right)^2\dfrac{d^2y}{dx^2}+2x\left(1+x^2\right)\dfrac{dy}{dx}=$ \underline{\hspace{2cm}}
\qn{17} If $f(x)=\min\left\{|x|,|x-2|,2-|x-1|\right\}$, then the number of points at which $f(x)$ is non-differentiable is \underline{\hspace{2cm}}
\qn{18} If $f(x)=\begin{cases}ax+b & ;\ x\le-1\\ ax^3+x+2b & ;\ x>-1\end{cases}$ $f(x)$ is differentiable at $x=-1$ then $2a+3b$ is \underline{\hspace{2cm}}
\end{multicols}
\end{document}
